% LSPV-4 REPORT::FORKFACTOR1::2026-10-12T01:00:00Z
\documentclass[11pt]{article}
\usepackage[a4paper,margin=0.86in]{geometry}
\usepackage{amsmath,amssymb,amsthm,booktabs,array,enumitem,hyperref}
\hypersetup{colorlinks=true,urlcolor=blue,linkcolor=blue}
\newcommand{\EXACT}{\textbf{PROVED/EXACT}}
\newcommand{\KNOWN}{\textbf{KNOWN THEOREM}}
\newcommand{\COMP}{\textbf{COMPUTATIONAL}}
\newcommand{\COND}{\textbf{CONDITIONAL}}
\newcommand{\CAND}{\textbf{CANDIDATE}}
\newcommand{\OPEN}{\textbf{OPEN}}
\newcommand{\RETIRED}{\textbf{FALSIFIED/RETIRED}}
\newtheorem{theorem}{Theorem}
\newtheorem{lemma}{Lemma}
\title{LSPV-4 FORKFACTOR1: Replacement-Covariant Forks Factor Through the Source}
\author{LSPV-4 Universal-Core Bridge}
\date{2026-10-12 01:00 UTC}

\begin{document}
\maketitle

\begin{abstract}
\textbf{Status: OPEN.} P versus NP remains open.  This run composes two
verified corpus boundaries: replacement covariance makes an extracted
commitment independent of the SAT implementation, while a total
polynomially verifiable two-challenge opening relation is searchable from the
source under the counterfactual $P=NP$.  Therefore a wrapper-robust fork is a
source functional, not evidence that the original one-bit decider exposed or
paid for the fork.  If the fork instead depends on implementation traces, it
fails the replacement-covariance requirement needed for an arbitrary-decider
claim.  The unqualified intermediate claim
\texttt{WRAPPER-COVARIANT-FORK-EXPOSURE1} is \RETIRED.  The surviving target
is a total source compiler followed by a model-matched unconditional
forbidden endpoint; that complete arrow is \OPEN{} and is already
separation-strength.
\end{abstract}

\section{Preflight, predecessor, and obstacles}

The cyclic predecessor is LSPV-3.  Its newest report,
\path{research/agents/lspv-3/2026-10-12_0036_LSPV-3_ONETRANS1.tex}, was read
back from GitHub, is substantive, and contains both the expected
\texttt{ONETRANS1} header and final verification sentinel.  No recovery write
was required.

The two obstacles selected before substantive work were:
\begin{enumerate}[leftmargin=1.7em]
\item obtaining a same-commitment, different-challenge pair from a
deterministic one-bit interface without importing it through the extractor;
\item transferring any such pair to a resource forced in the same execution
of every uniform deterministic polynomial-time SAT decider.
\end{enumerate}
Primary literature was checked first.  Standard special soundness is an
interactive statement about two accepting conversations, not a static
one-bit decision interface \cite{Damgaard10}.  Bogdanov--Trevisan show that
generic nonadaptive worst-case/average-case reductions for NP face a
polynomial-hierarchy-collapse barrier \cite{BT06}; this does not itself prove
our bridge impossible, but blocks treating an LWE-like analogy as a theorem.
Baker--Gill--Solovay and Aaronson--Wigderson remain the applicable
relativization and algebrization warnings \cite{BGS75,AW09}.  The
Fortnow--Lipton--van Melkebeek--Viglas lower bound is for one ordinary machine
in a particular simultaneous time--space regime, not for a semantic fork or
an arbitrary implementation trace \cite{FLMV05}.

\section{Clean core, semantic compatibility, and target}

Use the current clean core
\[
 \mathcal H_{2,2}=\{F:F\text{ is duplicate-free proper Monotone }
 3\mathrm{SAT}\text{-}(2,2)\}.
\]
For $F=C_1\wedge\cdots\wedge C_m$ define
\[
 H(F)_{\alpha,j}=1\quad\Longleftrightarrow\quad \alpha\models C_j.
\]
The decision bit is whether $H(F)$ has an all-one row.  A canonical indexed
encoding of an $n$-variable member has length
$N=|\operatorname{enc}(F)|=\Theta(n\log n)$.  Neither the definition of
$H(F)$ nor the one-bit output requires a decider to materialize $H(F)$, a
witness, a transcript, or an opening.

Let $\mathsf{PolySAT}$ be the codes of total uniform deterministic
polynomial-time exact SAT deciders.  Under the contradiction hypothesis
$P=NP$, this class is nonempty; every member may have its own exponent.  The
unrestricted lower-bound target remains
\[
 \forall A\in\operatorname{Dec}_{\rm uniform}(3\mathrm{SAT})\;
 \forall d\;\forall N_0\;\exists F\in\mathcal H_{2,2}:
 |F|\ge N_0\ \wedge\ T_A(F)>|F|^d,                       \tag{1}
\]
which is \OPEN.

Fix a polynomially balanced polynomial-time predicate
$V(F,c,e,z)$.  Its two-challenge pair relation is
\[
 Q(F,c,e_0,z_0,e_1,z_1):=[e_0\ne e_1]\wedge
 V(F,c,e_0,z_0)\wedge V(F,c,e_1,z_1).                    \tag{2}
\]
For an implementation $A$, a proposed compiler may name a prechallenge
commitment $c_A(F)$.  The SAME\_REALIZATION ledger contains $F$, $A$, the
compiler, challenge generation, all oracle calls, $c$, both openings,
validation, extraction, output, elapsed time and peak workspace.

\section{Exact covariance--source factorization}

\begin{theorem}[\texttt{FORK-COVARIANCE-FACTORIZATION1} --- \EXACT]
Assume $\mathsf{PolySAT}\ne\varnothing$.  For every
$B\in\mathsf{PolySAT}$ let $\kappa_B$ be the constant program transformer
$\kappa_B(A)=B$.  If
\[
 c_{\kappa_B(A)}(F)=c_A(F)
 \quad\text{for all }A,B\in\mathsf{PolySAT},\ F\in\mathcal H_{2,2}, \tag{3}
\]
then there is a unique source map $r$ such that
$c_A(F)=r(F)$ for every $A,F$.
\end{theorem}

\begin{proof}
For arbitrary $A,B,F$, equation (3) and $\kappa_B(A)=B$ give
$c_B(F)=c_A(F)$.  Fix $A_0$ and put $r(F)=c_{A_0}(F)$.  Independence of the
choice of $A_0$ and uniqueness follow immediately.  This is the commitment
specialization of \texttt{REPLACEMENT-COVARIANCE-FACTORIZATION1}; it is not a
runtime lower bound.
\end{proof}

\begin{lemma}[Total source-pair search --- \COND]
If (2) is polynomially balanced, polynomial-time decidable, and nonempty for
every promised $F$, then $P=NP$ implies a uniform polynomial-time algorithm
that returns the lexicographically first tuple in $Q_F$.
\end{lemma}

\begin{proof}
Whether a given prefix extends to a tuple in $Q_F$ is in NP.  Under $P=NP$,
polynomially many prefix-extension decisions reconstruct the first tuple.
All tuple lengths and verifier calls are polynomial in $N$ by hypothesis.
\end{proof}

\begin{theorem}[\texttt{FORK-TRANSFER-DICHOTOMY1} --- \EXACT]
For an alleged all-decider fork bridge, exactly the following guard applies.
\begin{enumerate}[leftmargin=1.7em]
\item If the prechallenge commitment obeys full exact-replacement covariance,
it factors through $F$.  If (2) is also total, balanced and efficiently
verifiable, then under $P=NP$ the fork is source-searchable.  Its existence
does not prove that the original execution of $A$ exposed, stored, or paid for
it.
\item If the named commitment or opening obligation changes under an exact
constant replacement, then the object is implementation-sensitive and does
not cover every $A\in\mathsf{PolySAT}$ without an additional typed
orbit-coverage theorem.
\end{enumerate}
\end{theorem}

\begin{proof}
The first branch is the preceding theorem and lemma.  Failure of their
covariance premise is exactly the second branch.  The alternatives are
exhaustive.  Notice that source search constructs a new algorithm; charging
its work to the original call of $A$ would violate SAME\_REALIZATION.  If the
source fork is then mapped to an unconditionally forbidden ordinary-machine
realization, that complete map is the desired separation, not a previously
available intermediate theorem.
\end{proof}

Thus \texttt{WRAPPER-COVARIANT-FORK-EXPOSURE1} is \RETIRED{} in its
unqualified form.  This does not retire
\texttt{EXTERNAL-CHALLENGE-SOURCE-PASS1}; it classifies that target as a
source compiler and leaves its totality, extraction consequence, costs and
forbidden endpoint \OPEN.

\section{Bounded reproducible audit --- \COMP}

The dependency-free script
\path{experiments/review/fork_covariance_factorization_audit.py} completed in
0.268 seconds, with each invocation below the 30-second limit.  It exhaustively
checked all $2^{12}=4096$ Boolean commitment maps on four implementations and
three sources.  Exactly $2^3=8$ were implementation-invariant and every one
factored through the source; factorization failures were zero.

It also checked all $2^{16}=65{,}536$ Boolean verifier tables on two sources
and one-bit commitment, challenge and reply alphabets.  There were $42{,}849$
relations total on both sources and $105{,}984$ source instances admitting a
pair.  Lexicographic reconstruction had zero failures.  DIRECT, COPY, DEAD,
DELAY, INTERP and RECOMP controls kept source-only commitments fixed; a
trace-sensitive tag failed covariance on all three sources.  Named Horn,
2SAT, affine/Tseitin, bounded-treewidth, independent-product, padding and
alternate-encoding controls were recorded.  The artifact
\path{results/raw/review/2026-10-12_fork_covariance_factorization.json}
explicitly makes no asymptotic claim.

\section{Hostile controls and barrier audit}

Horn forward chaining, 2SAT strongly connected components, Gaussian
elimination for affine/Tseitin instances, and bounded-treewidth dynamic
programming show that rich compatibility descriptions can coexist with
polynomial semantic shortcuts.  Independent products permit componentwise
solution.  Padding and redundant clauses alter $N$ or apparent state volume
without forcing work; alternative encodings alter traces and syntax.
One-bit DIRECT/COPY/DEAD/DELAY/INTERP/RECOMP wrappers preserve the decision
language while disrupting any trace-specific opening.  None of these finite
controls extrapolates to (1).

The factorization proof relativizes: exact replacement programs can use the
same oracle, and equality still gives a source map.  The prefix-search lemma
also relativizes when the corresponding $P^O=NP^O$ premise is made.  The
argument is algebraic-oracle stable and supplies no nonalgebrizing hinge.
It is not a Natural-Proofs argument, but it also supplies no circuit lower
bound.  FLMV can be invoked only after the same ordinary machine is placed in
its stated time--space regime; a source fork is not such a placement.

\section{Post-quantum and cross-domain transfer ledger}

\textbf{EXACT OBJECT/THEOREM TRANSFER:} the special-soundness pair relation
is represented exactly by (2), and the replacement-covariance theorem applies
exactly to its prechallenge commitment.  Collision entropy applies exactly to
any declared finite opening map, but measures lossiness rather than arbitrary
decision time.

\textbf{ANALOGY/CANDIDATE needing a new bridge:} cryptographic collision
extraction or special soundness could help only after a source-specific
construction proves that two valid openings yield an explicit forbidden
algebraic witness.  An LWE-style external black-box extraction architecture
would need a worst-case source compiler and must survive known NP
worst-case/average-case barriers.

\textbf{FAILED TRANSFER:} Schur squares, quadratic-relation ranks, Hilbert
functions and syzygies do not transfer merely because a fork exists; no code
or graded algebra is forced by an arbitrary SAT decider.  Random-oracle
programming, computational binding and post-quantum hardness assumptions are
inadmissible in an unconditional P-versus-NP proof.  As always, any hidden
algebraic structure might be an algorithmic shortcut rather than a hardness
source.

In other domains, observational equivalence and quotient constructions are
useful structural analogies for source factorization, but no control,
distributed, VLSI or communication lower bound transfers to the centralized
explicit-input SAT model without a compiler.  This classification is
\textbf{useful structural analogy requiring a new transfer theorem}, not a
direct lower bound.

\section{CORPUS REVIEWED}

I surveyed the current repository tree and search history, then read root
\path{main.tex}, \path{AGENTS.md}, \path{CODEX.md}, the complete
\path{.codex/skills/pnp-research/SKILL.md}, its mode/protocol references, all
current task specifications surfaced under \path{tasks/}, and actual
\path{main4.tex}/\path{main5.tex}.  I inspected the newest verified MAIN
snapshot and synthesis, latest LSPV-1 through LSPV-4 reports, current
MODELBRIDGE-0016/0018/0019, governance, and relevant historical LSPV-4
PCLOSE1/RANDCLOSE1/REACQ1/ORDERBIND1 reports.  I also inspected the underlying
HEC, VSC--HEC, G-PPC--HEC, CRS--PPC and CCSPC causal explorations and TASK-007,
rather than relying only on summaries.  Repository search was used to find
earlier covariance, fork, special-soundness, reacquisition and source-pass
formulations and their cited experiment artifacts.

The actual ARC manuscript \path{main4.tex} and WRAPPER manuscript
\path{main5.tex} remain construction programs with open all-input acquisition
and favorable replay/source-pass premises.  Neither manuscript proves the
arbitrary-decider fork implication above, and this report does not silently
rename or strengthen their objects.

\section{DUPLICATION / CONTRADICTION AUDIT}

MODELBRIDGE-0016 already proves replacement-covariance factorization;
SYNTHESIS-0018 already proves the opening-acquisition trilemma and conditional
total-pair search; ONETRANS1 already proves deterministic fixed-replay
uniqueness.  The new contribution is their exact composition at the
prechallenge commitment, producing a two-branch fork-transfer dichotomy and
a bounded joint regression test.  It is not a new proof of either ingredient.
There is no contradiction with ORDERBIND1: its still-open semantic binding
target is now narrowed to a source-specific forbidden-witness theorem, not an
internal-exposure claim.  There is no contradiction with HEC/CRS--PPC: their
restricted causal models do not provide full replacement covariance or an
all-decider compiler.

\section{WHAT CHANGED RELATIVE TO ROOT MAIN}

Root MAIN already retires physical order as an arbitrary runtime obligation
and keeps semantic opening binding, externally bounded source pass,
arbitrary-decider acquisition and the forbidden realization open.  This run
adds one exact reconciliation: if semantic opening binding is required to be
invariant across all exact replacement implementations, the commitment is
source-only; total polynomially verifiable forks are then searchable from the
source under $P=NP$.  The root LSPV-4 region should therefore retire generic
wrapper-covariant fork exposure while preserving the source compiler and
endpoint as \OPEN.

\section{ADOPTED / REJECTED / RETIRED PRIOR COMPONENTS}

\textbf{Adopted:} replacement covariance, fixed-replay uniqueness, the
opening-acquisition trilemma, exact encoded-length accounting,
SAME\_REALIZATION, and ORDERBIND1's demand for a source-specific semantic
conflict.  \textbf{Rejected as insufficient:} ordinary witness multiplicity,
trace inspection, latent semantic-state counts, and imported external
challenges without totality/cost proofs.  \textbf{FALSIFIED/RETIRED:}
\texttt{WRAPPER-COVARIANT-FORK-EXPOSURE1}, meaning the claim that a fork both
survives every exact replacement and witnesses implementation-specific
exposure by the original one-bit decider.  \textbf{Preserved OPEN:}
\texttt{EXTERNAL-CHALLENGE-SOURCE-PASS1}, a construction-specific forbidden
endpoint, typed restricted-orbit coverage if full replacement is deliberately
excluded, nonrelativization, nonalgebrization, and (1).

\section{STALE OR OVERSTATED CLAIMS FOUND}

Any statement that a wrapper-invariant extracted object is nevertheless an
implementation-sensitive resource is overstated when the wrapper family
contains all constant exact replacements.  Any statement that a total NP
fork relation is itself a hard resource under the $P=NP$ counterfactual is
also stale: prefix search makes it polynomial.  Any attribution of that
search cost to the original decider violates SAME\_REALIZATION.  Root MAIN
did not assert those strengthened readings; the ownership-region update
makes the guard explicit.

\section{Status and precise handoff}

\begin{center}\small
\begin{tabular}{p{0.48\linewidth}p{0.19\linewidth}p{0.25\linewidth}}
\toprule
Claim & Status & Scope\\ \midrule
\texttt{FORK-COVARIANCE-FACTORIZATION1} & \EXACT & Full constant replacement family\\
Total source-pair search under $P=NP$ & \COND & Balanced polynomial verifier\\
Finite joint audit & \COMP & 69,632 exhaustive objects; zero failures\\
\texttt{WRAPPER-COVARIANT-FORK-EXPOSURE1} & \RETIRED & Unqualified internal-exposure route\\
\texttt{EXTERNAL-CHALLENGE-SOURCE-PASS1} & \OPEN & Total source compiler and costs\\
Forbidden endpoint / (1) / P versus NP & \OPEN & Separation-strength\\
\bottomrule
\end{tabular}
\end{center}

The next task is not to extract another transcript from an arbitrary decider.
It is to give one explicit source-only relation $Q$ on the full canonical
clean core and prove: (i) total common-commitment forks, (ii) two openings
yield a construction-specific forbidden witness, and (iii) the resulting
ordinary machine lies in an unconditional model-matched forbidden regime
with all costs charged.  Failure at any arrow leaves P versus NP \OPEN.

\begin{thebibliography}{9}
\bibitem{Damgaard10} I. Damg\aa rd, \emph{On $\Sigma$-protocols}, lecture
notes, 2010, \url{https://www.cs.au.dk/~ivan/Sigma.pdf}.
\bibitem{BT06} A. Bogdanov and L. Trevisan, ``On worst-case to average-case
reductions for NP problems,'' \emph{SIAM J. Comput.} 36(4), 2006,
\url{https://lucatrevisan.github.io/pubs/redux-sicomp.pdf}.
\bibitem{BGS75} T. Baker, J. Gill and R. Solovay, ``Relativizations of the
$P=?NP$ question,'' \emph{SIAM J. Comput.} 4(4), 1975,
doi:10.1137/0204037.
\bibitem{AW09} S. Aaronson and A. Wigderson, ``Algebrization: a new barrier
in complexity theory,'' \emph{ACM Trans. Comput. Theory} 1(1), 2009.
\bibitem{FLMV05} L. Fortnow, R. Lipton, D. van Melkebeek and A. Viglas,
``Time-space lower bounds for satisfiability,'' \emph{J. ACM} 52(6), 2005.
\end{thebibliography}

% END_OF_VERIFIED_REPORT::LSPV-4::FORKFACTOR1
\end{document}
